% !TEX program = xelatex
% Lernplatt - by Can Siebert
% Kurs 22 (Bo 143) - Tag 6 - Lehrskript
% Omnichannel-Strategien & Operationalisierung digitaler Vertriebsstrategien
%
% Self-contained xelatex source. Kein biblatex, keine externen Bilder.

\documentclass[11pt,a4paper,oneside]{article}

% --- Encoding & Sprache --------------------------------------------------
\usepackage{polyglossia}
\setdefaultlanguage{german}

% --- Geometrie -----------------------------------------------------------
\usepackage[a4paper,top=2cm,bottom=2cm,outer=2.5cm,inner=2.5cm,bindingoffset=1.5cm]{geometry}

% --- Fonts ---------------------------------------------------------------
\usepackage{fontspec}
\defaultfontfeatures{Ligatures=TeX}

\IfFontExistsTF{Charter}{%
  \setmainfont{Charter}[Numbers=OldStyle]%
}{%
  \IfFontExistsTF{Noto Serif}{%
    \setmainfont{Noto Serif}%
  }{%
    \setmainfont{Liberation Serif}%
  }%
}

\IfFontExistsTF{Source Sans 3}{%
  \setsansfont{Source Sans 3}%
}{%
  \IfFontExistsTF{Source Sans Pro}{%
    \setsansfont{Source Sans Pro}%
  }{%
    \setsansfont{Liberation Sans}%
  }%
}

\newcommand{\caveatfontfallback}{\itshape}
\IfFontExistsTF{Caveat}{%
  \newfontfamily\caveatfont{Caveat}[Scale=1.15]%
}{%
  \newcommand\caveatfont{\caveatfontfallback}%
}

\setmonofont{Liberation Mono}

% --- Mikro-Typografie ----------------------------------------------------
\usepackage{microtype}
\linespread{1.15}

% --- Farben & Boxen ------------------------------------------------------
\usepackage{xcolor}
\definecolor{lpInk}{RGB}{20,28,38}
\definecolor{lpAccent}{RGB}{22,90,140}
\definecolor{lpAccentLight}{RGB}{210,228,240}
\definecolor{lpAmber}{RGB}{222,138,30}
\definecolor{lpAmberLight}{RGB}{255,243,217}
\definecolor{lpAmberBorder}{RGB}{196,118,18}
\definecolor{lpGray}{RGB}{96,104,114}
\definecolor{lpGrayLight}{RGB}{238,240,243}

\usepackage[most]{tcolorbox}
\tcbuselibrary{breakable,skins}

\newtcolorbox{eselsbox}[1][Eselsbrücke]{%
  enhanced, breakable,
  colback=lpAmberLight,
  colframe=lpAmberBorder,
  boxrule=0.6pt,
  arc=2pt,
  borderline={0.4pt}{0pt}{lpAmberBorder, dashed},
  left=10pt,right=10pt,top=8pt,bottom=8pt,
  fonttitle=\sffamily\bfseries\color{lpAmberBorder},
  coltitle=lpAmberBorder,
  title={#1},
  attach boxed title to top left={xshift=8pt,yshift=-8pt},
  boxed title style={size=small, colback=lpAmberLight, colframe=lpAmberBorder, boxrule=0.4pt, arc=1pt},
  top=14pt
}

\newtcolorbox{merkbox}[1][Merksatz]{%
  enhanced, breakable,
  colback=lpAccentLight,
  colframe=lpAccent,
  boxrule=0.6pt,
  arc=2pt,
  left=10pt,right=10pt,top=8pt,bottom=8pt,
  fonttitle=\sffamily\bfseries\color{white},
  coltitle=white,
  title={#1},
  attach boxed title to top left={xshift=8pt,yshift=-8pt},
  boxed title style={size=small, colback=lpAccent, colframe=lpAccent, boxrule=0.4pt, arc=1pt},
  top=14pt
}

% --- Listen --------------------------------------------------------------
\usepackage{enumitem}
\setlist{itemsep=2pt, topsep=4pt, parsep=0pt}
\setlist[itemize,1]{label={\textbullet},leftmargin=1.6em}
\setlist[itemize,2]{label={\textendash},leftmargin=1.4em}
\setlist[enumerate,1]{label=\arabic*., leftmargin=1.8em}

% --- Tabellen ------------------------------------------------------------
\usepackage{tabularx}
\usepackage{array}
\usepackage{booktabs}
\usepackage{longtable}

% --- Kopf-/Fusszeilen ----------------------------------------------------
\usepackage{fancyhdr}
\usepackage{lastpage}
\pagestyle{fancy}
\fancyhf{}
\renewcommand{\headrulewidth}{0.3pt}
\renewcommand{\footrulewidth}{0pt}
\fancyhead[L]{\footnotesize\sffamily\color{lpGray}Lernplatt \textperiodcentered{} Kurs 22 \textperiodcentered{} Tag 6}
\fancyhead[R]{\footnotesize\sffamily\color{lpGray}Omnichannel \& Operationalisierung}
\fancyfoot[L]{\footnotesize\sffamily\color{lpGray}\textcopyright{} Can Siebert}
\fancyfoot[R]{\footnotesize\sffamily\color{lpGray}\thepage{} / \pageref{LastPage}}

\fancypagestyle{plain}{%
  \fancyhf{}%
  \renewcommand{\headrulewidth}{0pt}%
  \fancyfoot[C]{\footnotesize\sffamily\color{lpGray}\thepage{} / \pageref{LastPage}}%
}

% --- Section-Styling -----------------------------------------------------
\usepackage[explicit]{titlesec}

\titleformat{\section}[block]
  {\sffamily\Large\bfseries\color{lpAccent}}
  {\thesection.}{0.6em}{#1}
  [{\vspace{2pt}\color{lpAccent}\titlerule[0.6pt]}]

\titleformat{\subsection}[block]
  {\sffamily\large\bfseries\color{lpInk}}
  {\thesubsection}{0.6em}{#1}

\titleformat{\subsubsection}[block]
  {\sffamily\normalsize\bfseries\color{lpInk}}
  {}{0pt}{#1}

\titlespacing*{\section}{0pt}{14pt}{8pt}
\titlespacing*{\subsection}{0pt}{10pt}{4pt}
\titlespacing*{\subsubsection}{0pt}{8pt}{2pt}

% --- Hyperlinks ----------------------------------------------------------
\usepackage[hidelinks,unicode]{hyperref}

% --- TOC -----------------------------------------------------------------
\renewcommand{\contentsname}{\sffamily\Large\bfseries\color{lpAccent}Inhalt}

% --- Helfer --------------------------------------------------------------
\newcommand{\akzent}[1]{{\caveatfont\color{lpAmber}#1}}
\newcommand{\fachbegriff}[1]{\textbf{#1}}
\newcommand{\quelle}[1]{{\footnotesize\color{lpGray}(#1)}}

% =========================================================================
\begin{document}

% --- COVER ---------------------------------------------------------------
\begin{titlepage}
  \pagestyle{empty}
  \vspace*{1.5cm}
  \noindent{\sffamily\footnotesize\color{lpGray}LERNPLATT \textperiodcentered{} BY CAN SIEBERT}\\[2pt]
  \noindent{\color{lpAccent}\rule{\linewidth}{1.2pt}}\\[4pt]
  \noindent{\sffamily\footnotesize\color{lpGray}MODUL 3 \textperiodcentered{} TAG 6 \textperiodcentered{} KURS 22 (Bo 143) \textperiodcentered{} 5.\,Juni 2026}

  \vspace{3cm}
  {\sffamily\fontsize{30}{34}\selectfont\bfseries\color{lpInk}Omnichannel\\\& Operationalisierung\par}

  \vspace{0.7cm}
  {\sffamily\large\color{lpGray}Vom Kanalmix zum digitalen Vertriebs­betriebssystem --\\Prozesse, Daten, Rollen, KPIs.\par}

  \vspace{1.5cm}
  {\caveatfont\LARGE\color{lpAmber}Lehrskript -- Tag 6 von 16}\par

  \vfill

  \noindent\begin{minipage}{0.55\linewidth}
    {\sffamily\footnotesize\color{lpGray}Lernpfad:}\\
    {\sffamily\small\color{lpInk}Wissen \textrightarrow{} Anwendung \textrightarrow{} Management}\\[10pt]
    {\sffamily\footnotesize\color{lpGray}Bearbeitungsdauer:}\\
    {\sffamily\small\color{lpInk}1 Kurstag}
  \end{minipage}\hfill
  \begin{minipage}{0.4\linewidth}
    \raggedleft
    {\sffamily\footnotesize\color{lpGray}Dozent:}\\
    {\sffamily\small\color{lpInk}Can Siebert}\\[10pt]
    {\sffamily\footnotesize\color{lpGray}Format:}\\
    {\sffamily\small\color{lpInk}Theorie \textperiodcentered{} Fallstudie \textperiodcentered{} Coaching}
  \end{minipage}

  \vspace{0.5cm}
  \noindent{\color{lpAccent}\rule{\linewidth}{0.6pt}}
\end{titlepage}

% --- LERNZIELE -----------------------------------------------------------
\section*{Lernziele dieses Tages}
\addcontentsline{toc}{section}{Lernziele dieses Tages}
\thispagestyle{plain}

\noindent An Tag~5 haben Sie eine kundenzentrierte Kanalstrategie abgeleitet. Heute geht es um die Umsetzungs- und Steuerungslogik: \emph{Wie wird aus einer Kanalstrategie ein steuerbares Omnichannel-System -- mit klaren Prozessen, sauberen Daten, definierten Rollen und gemeinsamen KPIs?}

\subsection*{L1 -- Wissen (Reproduktion und Einordnung)}
\begin{itemize}
  \item Omnichannel-Strategien verstehen und von Multichannel-Ansätzen unterscheiden.
  \item Zusammenspiel von Online- und Offline-Kanälen einordnen.
  \item Grundlagen der Operationalisierung digitaler Vertriebsstrategien kennen.
  \item Rollen von Vertrieb, Marketing, IT, Kundenservice und Management verstehen.
  \item Relevante KPIs für digitale Vertriebssteuerung benennen.
\end{itemize}

\subsection*{L2 -- Anwendung (Transfer auf konkrete Situationen)}
\begin{itemize}
  \item Omnichannel-Ansätze auf konkrete Vertriebssituationen übertragen.
  \item Kanalübergänge entlang der Customer Journey gestalten.
  \item Prozesse, Verantwortlichkeiten und Datenflüsse strukturieren.
  \item Digitale Vertriebsmaßnahmen priorisieren und in eine Roadmap überführen.
\end{itemize}

\subsection*{L3 -- Management (Entscheidung und Verantwortung)}
\begin{itemize}
  \item Omnichannel-Strategie als Führungs- und Architekturentscheidung bewerten.
  \item Zielkonflikte zwischen Kundenerlebnis, Vertriebseffizienz, Kosten, Datenqualität und Organisationsfähigkeit managen.
  \item Entscheidungen unter Unsicherheit, Budgetrestriktionen und Zeitdruck treffen.
  \item Verantwortung für eine skalierbare digitale Vertriebsarchitektur übernehmen.
\end{itemize}

\begin{merkbox}[Roter Faden]
Omnichannel bedeutet nicht ``\emph{überall sein}''. Es bedeutet, dass der Kunde an jedem Touchpoint dasselbe Unternehmen erlebt -- mit konsistenten Informationen, klaren Zuständigkeiten und ohne Medienbrüche. Wer das nicht operationalisiert, hat einen Kanalmix, kein Vertriebsbetriebssystem.
\end{merkbox}

\clearpage

% --- 1. EINSTIEG ---------------------------------------------------------
\section{Einstieg: Vom Kanal zum Betriebssystem}

In vielen Vertriebsorganisationen ist die folgende Situation Alltag: Es gibt eine Website, ein paar LinkedIn-Aktivitäten, ein CRM, einen Newsletter, einen Außendienst -- aber sie arbeiten kaum zusammen. Marketing misst Klicks, Vertrieb misst Abschlüsse, Kundenservice dokumentiert in Excel, Außendienst kennt die digitalen Kampagnen oft nicht. Das Ergebnis: Kunden erleben Brüche, Leads verfallen, das Management entscheidet auf Basis unvollständiger Daten. \quelle{vgl.\ Verhoef et al., 2015; Brynjolfsson et al., 2013}

\subsection{Vier typische Lücken in der Umsetzung}

\begin{enumerate}
  \item \fachbegriff{Strategie ohne operative Übersetzung.} Eine Kanalentscheidung ist getroffen -- aber niemand weiß genau, wer was bis wann tut.
  \item \fachbegriff{Prozesse ohne Datenbasis.} Es gibt Prozesse für Leadbearbeitung -- aber das CRM ist nicht verbindlich gepflegt, also fehlt die Steuerungsgrundlage.
  \item \fachbegriff{KPIs ohne gemeinsame Sicht.} Marketing und Vertrieb berichten unterschiedliche Größen. Wer welche Verantwortung trägt, ist nicht klar.
  \item \fachbegriff{Maßnahmen ohne Roadmap.} Initiativen werden parallel gestartet, ohne Priorisierung -- die Organisation wird überfordert, der Erfolg verteilt sich.
\end{enumerate}

\subsection{Was dieser Tag konkret leistet}

Wir bauen das Omnichannel-Modell systematisch auf: Begriffsklärung (Sektion 2), Kanalübergänge entlang der Journey (Sektion 3), Operationalisierung in Prozesse und Rollen (Sektionen 4--5), KPI-Architektur (Sektion 6), Roadmap-Logik (Sektion 7) und Fallstudie ``IndustrialPro Services'' (Sektion 8). Sektion 9 behandelt die Management-Perspektive.

\begin{eselsbox}[Eselsbrücke: \akzent{S-P-D-R} -- vier Bausteine der Operationalisierung]
Vier Buchstaben, vier Bausteine jeder Umsetzung:\\[2pt]
\textbf{S}trategie -- klar formuliertes Zielbild.\\
\textbf{P}rozesse -- definierter Ablauf für Lead, Übergabe, Nachverfolgung.\\
\textbf{D}aten -- verbindliches CRM, klare Datenflüsse.\\
\textbf{R}ollen -- wer entscheidet, wer berichtet, wer setzt um.\\[2pt]
\emph{Wer ``S-P-D-R'' nicht beantworten kann, hat eine Initiative -- keine Operationalisierung.}
\end{eselsbox}

% --- 2. SINGLE / MULTI / CROSS / OMNI ------------------------------------
\section{Single, Multi, Cross, Omni -- vier Stufen}

Die Begriffe werden oft unscharf verwendet. Eine saubere Abgrenzung ist Voraussetzung, um zu bewerten, wo eine Organisation steht und wohin sie sich bewegen will. \quelle{Verhoef et al., 2015; Rigby, 2011}

\subsection{Vier Stufen der Kanalintegration}

\begin{enumerate}
  \item \fachbegriff{Single Channel.} Ein Kanal, eine Logik. Klassisch: nur Außendienst, nur Webshop, nur Telefon. Robust, aber begrenzt.
  \item \fachbegriff{Multichannel.} Mehrere Kanäle nebeneinander, jeweils isoliert betrieben. Webshop \emph{und} Außendienst, aber ohne Datenaustausch. Häufigster Zustand mittelständischer Organisationen.
  \item \fachbegriff{Cross-Channel.} Kanäle sind teilweise verbunden -- Kunden können Kanäle wechseln, sehen aber noch Brüche. Z.\,B.\ Online-Kauf, aber Reklamation nur in der Filiale.
  \item \fachbegriff{Omnichannel.} Kanäle sind \emph{integriert}, der Kunde erlebt ein einheitliches Unternehmen. Konsistente Informationen, Preise, Bestände, Daten, Ansprechpartner -- unabhängig vom Kanal.
\end{enumerate}

\subsection{Was ``integriert'' konkret bedeutet}

Omnichannel ist mehr als ``alle Kanäle haben das gleiche Logo''. Integration bedeutet:

\begin{itemize}
  \item Der Kunde sieht überall konsistente Preise, Verfügbarkeiten und Lieferzeiten.
  \item Der Anbieter sieht über alle Kanäle einen einheitlichen Kundendatensatz -- nicht fünf Insellösungen.
  \item Eine Anfrage im Webchat ist im CRM sichtbar, wenn der Außendienst zum Termin geht.
  \item Marketing-Kampagnen sind dem Vertrieb bekannt, bevor der Kunde anruft.
  \item KPIs werden über Kanäle aggregiert betrachtet, nicht nur pro Kanal.
\end{itemize}

\begin{merkbox}[Omnichannel-Test]
Drei Fragen, die ein Omnichannel-System bestehen muss:
\begin{enumerate}
  \item Erlebt der Kunde einen konsistenten Anbieter -- unabhängig vom Kanal?
  \item Sehen wir intern \emph{einen} Kunden -- nicht fünf getrennte Datensätze?
  \item Können Marketing, Vertrieb, Service und Management auf \emph{eine} Datenbasis blicken, wenn sie eine Entscheidung treffen?
\end{enumerate}
Wenn ein ``Nein'' dabei ist, sind wir noch Multi- oder Cross-Channel, nicht Omni.
\end{merkbox}

% --- 3. KANALUEBERGAENGE -------------------------------------------------
\section{Kanalübergänge entlang der Customer Journey}

Omnichannel wird im Alltag spürbar an den \fachbegriff{Übergängen}: vom Online-Lead zum Beratungsgespräch, vom Webinar zur Angebotsanfrage, vom Webshop zum Außendienst, vom Kundenportal zum Servicekontakt. Genau dort entstehen die meisten Brüche -- und dort liegen die größten Hebel.

\subsection{Typische Übergänge}

\begin{itemize}
  \item \textbf{Lead-Übergänge.} Online-Anfrage \textrightarrow{} Vertrieb (Kontaktaufnahme innerhalb definierter Zeit).
  \item \textbf{Beratungs-Übergänge.} Anfrage mit Beratungsbedarf \textrightarrow{} Außendienst oder spezialisierter Innendienst.
  \item \textbf{Transaktions-Übergänge.} Angebot \textrightarrow{} Vertrag \textrightarrow{} Lieferung \textrightarrow{} Rechnung -- mit klarer Dokumentation im CRM.
  \item \textbf{Service-Übergänge.} Bestandskunde mit Anliegen \textrightarrow{} Kundenservice, optional zurück zu Vertrieb (Upsell).
  \item \textbf{Nachbestell-Übergänge.} Service \textrightarrow{} Vertrieb oder Self-Service-Portal, sobald Bedarf entsteht.
\end{itemize}

\subsection{Typische Brüche in der Praxis}

\begin{itemize}
  \item \textbf{Doppelte Datenerfassung.} Eine Anfrage wird in Web-Formular, E-Mail und CRM mehrfach erfasst -- mit Abweichungen.
  \item \textbf{Unklare Zuständigkeiten.} Nach dem Webinar weiß niemand, wer welchen Teilnehmer kontaktiert.
  \item \textbf{Widersprüchliche Informationen.} Website nennt Lieferzeit ``2 Wochen'', Außendienst sagt ``4 Wochen''.
  \item \textbf{Medienbrüche.} Ein Vorgang springt zwischen E-Mail, CRM, Excel-Tabelle, Outlook-Kalender -- ohne dass irgendwo der vollständige Stand sichtbar wird.
  \item \textbf{Fehlende Lead-Übergabe.} Marketing erzeugt einen Lead, aber niemand übernimmt ihn innerhalb von 24--48 Stunden -- der Kunde wendet sich an den Wettbewerber.
\end{itemize}

\quelle{vgl.\ Kannan \& Li, 2017}

\subsection{Gestaltung guter Übergänge}

Ein Übergang ist ``gut'' gestaltet, wenn drei Bedingungen erfüllt sind:

\begin{enumerate}
  \item \textbf{Trigger ist eindeutig.} Welcher Ereignistyp löst die Übergabe aus?
  \item \textbf{Empfänger ist klar.} Wer übernimmt -- nicht eine Abteilung, sondern eine konkrete Rolle.
  \item \textbf{Verbindlichkeit existiert.} Reaktionszeit ist definiert (z.\,B.\ ``erstkontakt innerhalb von 4 Arbeitsstunden'') und wird gemessen.
\end{enumerate}

\begin{eselsbox}[Eselsbrücke: \akzent{T-E-V} -- drei Säulen jedes Übergangs]
\textbf{T}rigger -- was löst aus?\\
\textbf{E}mpfänger -- wer übernimmt?\\
\textbf{V}erbindlichkeit -- bis wann, mit welcher Qualität?\\[2pt]
\emph{Wenn ein Übergang nicht alle drei beantwortet, fällt er ins Loch zwischen den Abteilungen.}
\end{eselsbox}

% --- 4. PROZESSE FUER DEN DIGITALEN LEAD ---------------------------------
\section{Prozessarchitektur für den digitalen Lead}

Das wichtigste Einzelobjekt im digitalen Vertrieb ist der \fachbegriff{digitale Lead}: eine Interessensbekundung, die digital entsteht -- über Formular, Download, Webinar-Anmeldung, Chat, Demo-Anfrage. Wer den Umgang mit diesem Objekt nicht beherrscht, hat keinen funktionierenden digitalen Vertrieb. \quelle{Ellis \& Brown, 2017; Homburg, 2020}

\subsection{Der Lead-Lifecycle in sechs Schritten}

\begin{enumerate}
  \item \fachbegriff{Aufnahme.} Lead entsteht über digitalen Touchpoint, wird automatisch ins CRM geschrieben.
  \item \fachbegriff{Qualifizierung.} Lead Scoring oder manuelle Bewertung: Wie wahrscheinlich ist Kaufabsicht, Budget, Entscheidungsbefugnis?
  \item \fachbegriff{Übergabe.} Qualifizierter Lead geht an die richtige Rolle (Innendienst, Außendienst, Service).
  \item \fachbegriff{Nachverfolgung.} Erstkontakt innerhalb definierter Zeit, weitere Schritte dokumentiert.
  \item \fachbegriff{Abschluss oder Verlust.} Lead wird zum Kunden -- oder als verloren markiert (mit Grund).
  \item \fachbegriff{Auswertung.} Pro Kanal, Kampagne, Zielgruppe: wie hoch waren Conversion, Abschlussquote, Customer Acquisition Cost?
\end{enumerate}

\subsection{Lead Scoring -- knapp erklärt}

\fachbegriff{Lead Scoring} ist die strukturierte Bewertung der Lead-Qualität anhand klar definierter Kriterien. Typisch:

\begin{itemize}
  \item \textbf{Demografie / Firmografie.} Passt die Person/Firma in die Zielgruppe? (Branche, Größe, Rolle).
  \item \textbf{Verhalten.} Welche Interaktionen hat sie ausgelöst? (Mehrfache Website-Besuche, Whitepaper-Download, Webinar-Teilnahme).
  \item \textbf{Bedarfssignale.} Konkrete Hinweise auf Kaufabsicht? (Demo-Anfrage, Preisseite besucht, Budget genannt).
  \item \textbf{Zeit.} Wie aktuell sind die Signale?
\end{itemize}

Aus diesen Dimensionen wird ein Score gebildet (z.\,B.\ 0--100), der eine Übergabeschwelle hat (z.\,B.\ ``ab 60 Punkten geht der Lead an den Vertrieb''). \quelle{vgl.\ Ellis \& Brown, 2017}

\subsection{Reaktionszeiten -- der unterschätzte Hebel}

Empirisch ist gut belegt: die ersten Stunden entscheiden. Ein Lead, der innerhalb von 4 Stunden kontaktiert wird, hat eine deutlich höhere Conversion-Wahrscheinlichkeit als einer, der nach drei Tagen kontaktiert wird. Reaktionszeit ist daher kein operativer Detailparameter, sondern eine \emph{strategische} Größe.

\begin{merkbox}[Lead-Prozess-Checkliste]
\begin{itemize}
  \item Jeder digitale Lead landet automatisch im CRM (kein E-Mail-Posteingang als Sammelstelle).
  \item Jeder Lead wird innerhalb definierter Zeit qualifiziert (Lead Scoring oder manuelle Triage).
  \item Jeder qualifizierte Lead hat einen Empfänger und eine maximale Reaktionszeit.
  \item Jeder verlorene Lead hat einen Verlustgrund -- dokumentiert für die Auswertung.
\end{itemize}
\end{merkbox}

% --- 5. ROLLEN, GOVERNANCE, ZUSAMMENARBEIT -------------------------------
\section{Rollen, Governance, Zusammenarbeit}

Omnichannel ist auch eine Organisationsfrage. Wenn Marketing, Vertrieb, Kundenservice und IT mit unterschiedlichen Zielen, KPIs und Tools arbeiten, hilft auch das beste CRM nicht. \quelle{Homburg, 2020}

\subsection{Fünf Rollen, die zusammenspielen müssen}

\begin{itemize}
  \item \fachbegriff{Marketing.} Zuständig für Reichweite, Kampagnen, Content -- aber \emph{nicht} für Umsatz allein.
  \item \fachbegriff{Vertrieb.} Zuständig für Abschluss, Kundenbeziehung -- aber \emph{nicht} losgelöst von der Leadqualität, die Marketing erzeugt.
  \item \fachbegriff{Kundenservice.} Zuständig für Bestandskundenpflege und Reklamationen -- aber auch ein wichtiger Datenlieferant für Vertrieb und Marketing.
  \item \fachbegriff{IT.} Zuständig für die technische Infrastruktur (CRM, Marketing Automation, Webshop, Schnittstellen) -- aber \emph{nicht} der Owner der Vertriebslogik.
  \item \fachbegriff{Management.} Zuständig für Zielsystem, Governance, Investitionsentscheidungen -- mit Verantwortung für die Architektur als Ganzes.
\end{itemize}

\subsection{Gemeinsame Erfolgs-KPIs als Brücke}

Ein bewährter Hebel: Marketing und Vertrieb teilen sich definierte gemeinsame KPIs (z.\,B.\ ``Anzahl qualifizierter Leads pro Monat'' und ``Conversion von Lead zu Angebot'' und ``Conversion von Angebot zu Abschluss''). Niemand kann allein erfolgreich sein -- beide tragen.

\subsection{Governance-Themen, die geklärt sein müssen}

\begin{itemize}
  \item Wer ist Eigentümer des CRM-Datenmodells?
  \item Welche Felder sind Pflicht, welche optional?
  \item Wer definiert Lead Scoring und Übergabeschwelle?
  \item Wer reagiert auf Beschwerden über Kanalübergänge?
  \item Welche Entscheidungen darf Marketing alleine treffen, welche Vertrieb, welche nur das Management?
\end{itemize}

\begin{merkbox}[Governance-Test für Omnichannel]
Wenn drei Personen aus drei Abteilungen die gleiche Frage zu einem Lead unterschiedlich beantworten, fehlt Governance. Wenn alle drei dasselbe sagen, funktioniert sie. Dieser Test ist primitiv -- und unbestechlich.
\end{merkbox}

% --- 6. KPI-ARCHITEKTUR -------------------------------------------------
\section{KPI-Architektur für digitale Vertriebssteuerung}

Eine Omnichannel-Architektur ist nur so gut wie ihre Steuerung. Ohne saubere KPIs entscheidet das Bauchgefühl -- und das ist im Mittelstand oft das Bauchgefühl jener Person, die am lautesten ihre Meinung vertritt. \quelle{Kaplan \& Norton, 1996; Ellis \& Brown, 2017}

\subsection{Drei Ebenen von KPIs}

\begin{enumerate}
  \item \textbf{Operative KPIs.} Tages- und Wochengrößen für die Linie -- Reaktionszeit, Anzahl Leads, Klickrate, Bearbeitungsstatus.
  \item \textbf{Taktische KPIs.} Monatliche Größen für die Bereichsleitung -- Conversion Rate, Kosten pro Lead, Abschlussquote, Anteil dokumentierter Leads.
  \item \textbf{Strategische KPIs.} Quartals-/Jahresgrößen für die Geschäftsleitung -- Customer Acquisition Cost, Customer Lifetime Value, Kanalbeitrag, Marge nach Kanal.
\end{enumerate}

\subsection{Kern-KPIs für digitale Vertriebssteuerung}

\begin{itemize}
  \item \fachbegriff{Anzahl qualifizierter Leads.} Nur Leads, die den Scoring-Schwellenwert erreichen.
  \item \fachbegriff{Reaktionszeit.} Zeit vom Lead-Eingang bis zum Erstkontakt durch Vertrieb.
  \item \fachbegriff{Conversion Rate.} Anteil Leads, die zur nächsten Stufe wechseln (Anfrage \textrightarrow{} Angebot \textrightarrow{} Abschluss).
  \item \fachbegriff{Kosten pro Lead (CPL).} Wie viel kostet ein qualifizierter Lead -- gesamt und pro Kanal?
  \item \fachbegriff{Customer Acquisition Cost (CAC).} Wie viel kostet ein \emph{gewonnener Kunde}?
  \item \fachbegriff{Customer Lifetime Value (CLV).} Wie viel Wert schafft ein Kunde über die gesamte Beziehung?
  \item \fachbegriff{Abschlussquote.} Anteil qualifizierter Leads, die zum Abschluss werden.
  \item \fachbegriff{Anteil vollständig dokumentierter Leads im CRM.} Steuerungsfähigkeit der Datenbasis.
  \item \fachbegriff{Kundenzufriedenheit nach digitalem Kontakt.} Verbindendes Maß für Service \& Vertrieb.
\end{itemize}

\subsection{Falle: KPI-Theater}

Eine häufige Falle: KPIs werden gemessen, aber niemand reagiert. Reports werden erstellt, aber nicht gelesen. \emph{KPI-Theater} ist gefährlich, weil es Steuerung simuliert, die nicht existiert. Eine sinnvolle Regel: Pro KPI gibt es eine Person, die entscheidet, was wann passiert, wenn der Wert kippt.

\begin{eselsbox}[Eselsbrücke: \akzent{L-R-C-K-K-A} -- sechs KPIs für die Linie]
Sechs Buchstaben, sechs operativ-taktische KPIs:\\[2pt]
\textbf{L}eads (qualifiziert)\\
\textbf{R}eaktionszeit\\
\textbf{C}onversion Rate\\
\textbf{K}osten pro Lead\\
\textbf{K}undenakquisitionskosten (CAC)\\
\textbf{A}bschlussquote\\[2pt]
\emph{Wer als Linienleitung diese sechs nicht wöchentlich auf dem Schreibtisch hat, steuert nicht.}
\end{eselsbox}

% --- 7. ROADMAP & PILOTLOGIK ---------------------------------------------
\section{Roadmap, Pilotlogik, Rollout}

Eine Omnichannel-Architektur entsteht nicht in einem ``Big Bang''. Sie wächst in Etappen -- über Pilot, Roll-out, Skalierung. Eine belastbare 12-Monats-Roadmap arbeitet typischerweise mit vier Quartalsschritten. \quelle{Rogers, 2016}

\subsection{Quartalslogik einer 12-Monats-Roadmap}

\begin{itemize}
  \item \textbf{Quartal 1 (Fundament).} CRM-Datenmodell, Lead-Felder, Rollen, gemeinsame KPIs definieren. Reaktionszeiten festlegen. Datenqualität bereinigen.
  \item \textbf{Quartal 2 (Pilot).} Lead Scoring einführen. Website-Anfrageprozesse verbessern. Schulung der beteiligten Rollen. Erste Übergaberegeln in einem Segment testen.
  \item \textbf{Quartal 3 (Roll-out).} Erweiterung auf weitere Segmente. Gemeinsames Dashboard einführen. E-Mail-Nurturing aufbauen. Erste Auswertung der Leadqualität.
  \item \textbf{Quartal 4 (Optimierung).} Komplexere Initiativen (z.\,B.\ Kundenportal-Pilot). Prozesse optimieren. Entscheidung über Skalierung der nächsten Stufe.
\end{itemize}

\subsection{Was \emph{nicht} ins Q1 gehört}

\begin{itemize}
  \item Neue Tool-Einführung, solange CRM-Datenqualität nicht steht.
  \item Großprojekte wie Kundenportal-Vollausbau -- besser als Pilot mit klar abgegrenztem Bestandskunden-Segment.
  \item Werbekampagnen mit großem Budget, solange Conversion-Pfade nicht stabil sind.
\end{itemize}

\subsection{Pilot statt Großprojekt}

Wo immer möglich, gilt: \emph{Pilot vor Roll-out}. Ein klar abgegrenzter Pilot (z.\,B.\ Kundenportal für 50 Bestandskunden mit Wartungsverträgen) liefert in 3--6 Monaten belastbare Erkenntnisse -- mit deutlich geringerem Risiko als ein 18-Monats-Großprojekt für alle Kunden.

\begin{merkbox}[Roadmap-Grundregel]
Bauen Sie das Fundament, bevor Sie die Fenster putzen. Tool-Einführungen, neue Kampagnen oder Portale wirken erst, wenn Datenbasis, Rollen und KPIs stehen. Die häufigste Roadmap-Falle ist nicht ein fehlendes Tool, sondern ein fehlendes Fundament.
\end{merkbox}

% --- 8. FALLSTUDIE-LEITFADEN ---------------------------------------------
\section{Fallstudie-Leitfaden: IndustrialPro Services}

In den Unterrichtseinheiten 4--5 bearbeiten Sie die Fallstudie ``IndustrialPro Services'' -- ein mittelständisches Unternehmen für Wartung, Reparatur und Modernisierung von Produktionsanlagen mit 38\,Mio.\,EUR Umsatz, 18 Vertriebs-, 3 Marketing- und 12 Service-Mitarbeitenden.

\subsection{Analytischer Aufbau}

\begin{enumerate}
  \item \textbf{Omnichannel-Schwächen identifizieren.} Welche Brüche existieren? Welche sind Daten-, welche Prozess-, welche Rollenprobleme?
  \item \textbf{Ideale Customer Journey für einen digitalen Service-Lead.} Vom Erkennen des Bedarfs bis zur Vertragsverlängerung -- mit konkreten Touchpoints.
  \item \textbf{Maßnahmen bewerten.} Welche der vorgeschlagenen Maßnahmen (CRM-Pflichtprozess, Lead Scoring, Website-Relaunch, E-Mail-Strecken, Kundenportal, Dashboard, Außendienst-Schulung) hat welche Priorität?
  \item \textbf{Priorisierung für 12 Monate.} Maximal 5 Maßnahmen in Quartalslogik.
  \item \textbf{Roadmap mit Verantwortlichkeiten.} Pro Maßnahme ein Eigentümer.
  \item \textbf{KPIs definieren.} 5--7 KPIs, die die digitale Vertriebsstrategie steuern.
  \item \textbf{Entscheidung unter Unsicherheit.} Welche attraktive Maßnahme wird trotz Verlockung \emph{nicht} sofort gewählt?
\end{enumerate}

\subsection{Erwartete Empfehlungsrichtung}

In einer typischen Musterlösung wird der Schwerpunkt auf \emph{CRM-Pflichtprozess}, \emph{Lead Scoring \& Übergaberegeln}, \emph{Website-Optimierung}, \emph{gemeinsames Dashboard} und \emph{Außendienst-Schulung} gelegt. Das Kundenportal wird \emph{als Pilot} für eine kleine Bestandskundengruppe begonnen, nicht als Großprojekt für alle Kunden. Die Entscheidung unter Unsicherheit: ``Trotz Attraktivität rollen wir das Kundenportal nicht breit aus, weil ohne stabile CRM-Prozesse und Rollen die zusätzliche Komplexität die Organisation überfordern würde.''

\begin{merkbox}[Argumentationsmuster für die Fallstudie]
Eine gute Omnichannel-Roadmap enthält:
\begin{enumerate}
  \item \textbf{Was zuerst.} Fundament-Maßnahmen (Daten, Rollen, KPIs) in Quartal 1.
  \item \textbf{Was als Pilot.} Komplexe Initiativen in begrenztem Segment.
  \item \textbf{Was zurückgestellt.} Großprojekte, die ohne Fundament Komplexität ohne Wert erzeugen würden.
  \item \textbf{Welche Verantwortung wo.} Pro Maßnahme ein Eigentümer, pro KPI ein Berichtender.
  \item \textbf{Welche Annahme wir prüfen.} Welche Frühindikatoren signalisieren, dass die Strategie kippt?
\end{enumerate}
\end{merkbox}

% --- 9. MANAGEMENT-PERSPEKTIVE ------------------------------------------
\section{Management-Perspektive: Zielkonflikte}

Auf Wissens- und Anwendungsebene sieht eine gute Omnichannel-Strategie wie eine sauber priorisierte Roadmap aus. Auf Managementebene geht es um Zielkonflikte, die nicht ``gelöst'', sondern \emph{entschieden} werden müssen -- vor allem darüber, welche Initiativen man bewusst nicht startet. \quelle{vgl.\ Brynjolfsson et al., 2013; Rogers, 2016}

\subsection{Fünf zentrale Zielkonflikte}

\begin{enumerate}
  \item \textbf{Kundenerlebnis vs.\ Vertriebseffizienz.} Mehr Touchpoints können das Kundenerlebnis verbessern -- bedeuten aber mehr Vertriebsaufwand pro Kunde.
  \item \textbf{Geschwindigkeit vs.\ Datenqualität.} Initiativen schnell starten ist möglich, aber wenn die Datenbasis nicht stimmt, multipliziert man Brüche.
  \item \textbf{Kosten vs.\ Organisationsfähigkeit.} Tools sind günstig, eine fähige Organisation ist teuer. Wer am Aufbau spart, bezahlt es später vielfach in Reibungsverlusten.
  \item \textbf{Initiativen-Anzahl vs.\ Umsetzungstiefe.} Drei Initiativen ordentlich umgesetzt schlagen sieben halb umgesetzte.
  \item \textbf{Marketing-Reichweite vs.\ Vertriebskapazität.} Mehr Leads erzeugen ist einfach -- mit ihnen sinnvoll umzugehen ist schwer. Wer den Vertrieb mit unqualifizierten Kontakten überflutet, verbrennt Kapazität.
\end{enumerate}

\subsection{Architekturentscheidung statt Tool-Auswahl}

Omnichannel ist \emph{keine} Tool-Entscheidung. Sie ist eine Architekturentscheidung über Kundenschnittstellen, Daten, Prozesse und Rollen. Sie wird auf Senior-Level entschieden -- typischerweise durch CEO, Chief Revenue Officer oder Vertriebsleitung gemeinsam mit IT. Wer Omnichannel an die IT delegiert, baut ein Datensystem; wer an Marketing delegiert, baut eine Kampagnenmaschine. Beides ist zu wenig.

\subsection{Senior-Level-Denken}

Aus den Coaching-Themen (UE 4) und dem Feedback (UE 7) ergibt sich ein einfacher Test: Eine senior-level-fähige Omnichannel-Entscheidung lässt sich in einem Satz erklären, der eine bewusste Verzichts­entscheidung enthält und das Frühwarnsignal benennt. Beispiel: ``Wir starten zuerst CRM-Pflichtprozess und Lead Scoring, das Kundenportal nur als Pilot mit 50 Bestandskunden -- wenn die CRM-Dokumentationsquote nach 3 Monaten nicht über 80\,\% steigt, stoppen wir das Portal-Projekt.''

\begin{eselsbox}[Eselsbrücke: \akzent{K-G-K-I-M}]
Fünf Zielkonflikte, die auf Senior-Level entschieden werden:\\[2pt]
\textbf{K}undenerlebnis vs.\ Vertriebseffizienz.\\
\textbf{G}eschwindigkeit vs.\ Datenqualität.\\
\textbf{K}osten vs.\ Organisationsfähigkeit.\\
\textbf{I}nitiativen-Anzahl vs.\ Umsetzungstiefe.\\
\textbf{M}arketing-Reichweite vs.\ Vertriebskapazität.\\[2pt]
\emph{Fünf Konflikte, fünf bewusste Entscheidungen -- keiner davon hat eine ``optimale Lösung''.}
\end{eselsbox}

\begin{merkbox}[Schluss-Merksatz für Tag 6]
Omnichannel ist nicht der Maximalausbau aller Kanäle. Es ist die Disziplin, an jedem Touchpoint das gleiche Unternehmen erlebbar zu machen -- mit konsistenten Informationen, klaren Verantwortlichkeiten und einer Datenbasis, auf die alle blicken. \emph{Wer das nicht operationalisiert, hat einen Werbeauftritt -- kein Vertriebsbetriebssystem.}
\end{merkbox}

% --- 10. LITERATUR -------------------------------------------------------
\clearpage
\section{Literatur}

Im Skript verwendete und für die Vertiefung empfohlene Quellen. Zitierstil: APA-7.

\begin{description}[style=nextline,leftmargin=18pt,labelindent=0pt,itemsep=4pt]
  \item[Brynjolfsson, E., Hu, Y.\,J., \& Rahman, M.\,S. (2013).] Competing in the age of omnichannel retailing. \emph{MIT Sloan Management Review, 54}(4), 23--29.
  \item[Ellis, S., \& Brown, M. (2017).] \emph{Hacking growth: How today's fastest-growing companies drive breakout success}. Crown Business.
  \item[Homburg, C. (2020).] \emph{Marketingmanagement: Strategie -- Instrumente -- Umsetzung -- Unternehmensführung} (7.\ Aufl.). Springer Gabler.
  \item[Kannan, P.\,K., \& Li, H. (2017).] Digital marketing: A framework, review and research agenda. \emph{International Journal of Research in Marketing, 34}(1), 22--45. \href{https://doi.org/10.1016/j.ijresmar.2016.11.006}{https://doi.org/10.1016/j.ijresmar.2016.11.006}
  \item[Kaplan, R.\,S., \& Norton, D.\,P. (1996).] \emph{The balanced scorecard: Translating strategy into action}. Harvard Business School Press.
  \item[Rigby, D. (2011).] The future of shopping. \emph{Harvard Business Review, 89}(12), 64--75.
  \item[Rogers, D.\,L. (2016).] \emph{The digital transformation playbook: Rethink your business for the digital age}. Columbia Business School Publishing.
  \item[Verhoef, P.\,C., Kannan, P.\,K., \& Inman, J.\,J. (2015).] From multi-channel retailing to omni-channel retailing. \emph{Journal of Retailing, 91}(2), 174--181. \href{https://doi.org/10.1016/j.jretai.2015.02.005}{https://doi.org/10.1016/j.jretai.2015.02.005}
\end{description}

% --- 11. GLOSSAR ---------------------------------------------------------
\clearpage
\section{Glossar}

Die wichtigsten Begriffe des Tages -- kurz definiert, in alphabetischer Reihenfolge.

\begin{description}[style=nextline,leftmargin=0pt,labelindent=0pt,itemsep=4pt]
  \item[Abschlussquote] Anteil qualifizierter Leads, der zum Vertragsabschluss führt.
  \item[CRM] Customer Relationship Management -- System zur strukturierten Erfassung und Steuerung von Kundenbeziehungen entlang des Lifecycles.
  \item[Cross-Channel] Vertriebsansatz mit teilweise verbundenen Kanälen, in denen Kunden Kanäle wechseln können -- aber Brüche erleben.
  \item[Customer Acquisition Cost (CAC)] Gesamtkosten, die anfallen, um einen neuen Kunden zu gewinnen.
  \item[Customer Lifetime Value (CLV)] Erwarteter wirtschaftlicher Wert eines Kunden über die gesamte Beziehung.
  \item[Customer Journey] Strukturierte Abfolge der Phasen, die ein Kunde durchläuft -- siehe Tag 5.
  \item[Dashboard] Visuelle Aggregation der wichtigsten KPIs für eine Zielgruppe (Linie, Bereichsleitung, Management).
  \item[Datenflüsse] Definierte Wege, auf denen Informationen zwischen Systemen, Rollen und Kanälen wandern.
  \item[Konsistenz (Omnichannel)] Eigenschaft, dass Kunden über alle Kanäle dieselben Informationen, Preise, Verfügbarkeiten und Ansprechpartner erleben.
  \item[KPI] Key Performance Indicator -- messbare Steuerungsgröße. Unterscheidung in operative, taktische und strategische KPIs.
  \item[Lead] Interessensbekundung einer Person oder Organisation, die in einen Vertriebsprozess überführt werden kann.
  \item[Lead Scoring] Strukturierte Bewertung der Lead-Qualität anhand definierter Kriterien (Demografie, Verhalten, Bedarfssignale, Aktualität).
  \item[Marketing Automation] Software-Logik zur automatisierten, regelbasierten Ansprache von Interessenten und Bestandskunden.
  \item[Medienbruch] Stelle in einem Prozess, an der ein Vorgang das Medium wechselt (z.\,B.\ E-Mail \textrightarrow{} Excel \textrightarrow{} CRM) -- typischer Verlustpunkt.
  \item[Multichannel] Vertriebsansatz mit mehreren parallel betriebenen, kaum integrierten Kanälen.
  \item[Omnichannel] Vertriebsansatz, in dem alle Kanäle integriert sind und Kunden überall ein konsistentes Unternehmen erleben.
  \item[Operationalisierung] Übersetzung einer Strategie in Prozesse, Rollen, Daten, Werkzeuge und steuerbare KPIs.
  \item[Reaktionszeit] Zeit vom Eingang eines Leads bis zum ersten Kontakt durch den Vertrieb.
  \item[Roadmap] Geordneter, etappierter Umsetzungsplan über mehrere Quartale; trennt Fundament, Pilot, Roll-out, Optimierung.
  \item[Rollen (Omnichannel)] Marketing, Vertrieb, Kundenservice, IT, Management -- mit klar definierten Verantwortlichkeiten.
  \item[Single Channel] Vertrieb über einen einzigen Kanal (z.\,B.\ nur Außendienst).
  \item[Touchpoint] Konkreter Berührungspunkt zwischen Kunde und Anbieter entlang der Journey.
  \item[Übergabe (Lead)] Strukturierte Weiterleitung eines qualifizierten Leads an die nächste Rolle, mit definierter Reaktionszeit.
  \item[Zielkonflikt (Omnichannel)] Spannungsfeld zwischen Kundenerlebnis, Geschwindigkeit, Kosten, Datenqualität, Organisationsfähigkeit und Vertriebskapazität.
\end{description}

\vspace{1cm}
\noindent{\color{lpAccent}\rule{\linewidth}{0.6pt}}\\[4pt]
{\footnotesize\sffamily\color{lpGray}Lernplatt \textperiodcentered{} by Can Siebert \textperiodcentered{} Kurs 22 (Bo 143) \textperiodcentered{} Tag 6 \textperiodcentered{} \textcopyright{} 2026}

\end{document}
